\chapter{Run, inspect, and reproduce the project}
\label{chap:reproduce}

\section{Three things that are easy to confuse}

The \textbf{source code} is the set of instructions in files ending in
\code{.jl}. The \textbf{results} are the numbers those instructions write
after a run. The \textbf{guide} is the document you are reading. Editing a
sentence in the guide does not run the model. Running the model does not
automatically turn every new result into a new PDF.

We therefore use separate steps: run the Julia code, prepare the document's
inputs, and compile the LaTeX document. This chapter explains the steps and
then walks through the scripts that connect them.

\section{Open the project folder in a terminal}

A \textbf{terminal} accepts text commands for the operating system. It is
different from the Julia prompt. The commands below belong in the terminal.
If you see the Julia prompt \code{julia>}, leave that session first or open
a separate terminal.

Extract the supplied project archive if you are starting from the archive.
Locate the folder named \code{moran_project}. Inside it, you should find
folders named \code{src}, \code{scripts}, \code{test}, \code{results},
and \code{guide}.

The command \code{cd} changes the terminal's current directory. The next
line uses a placeholder path. Replace \code{/path/to/moran_project} with
the actual path to your extracted folder. Keep the quotation marks if the
path contains spaces.

\begin{lstlisting}[language={},numbers=none]
cd "/path/to/moran_project"
\end{lstlisting}

To check where the terminal is looking, use these two commands:
\begin{lstlisting}[language={},numbers=none]
pwd
ls
\end{lstlisting}
The first prints the current directory. The second lists its contents. If
you do not see \code{scripts} and \code{src}, you are probably looking at
the wrong folder. This is a file-location problem, not a population-genetics
result.

\section{Run the small demonstration first}

Run this command from the project folder:
\begin{lstlisting}[language={},numbers=none]
julia --project=. scripts/demo.jl
\end{lstlisting}
The word \code{julia} starts Julia. The option \code{--project=.} selects
the current directory as the Julia project environment. A single dot means
``the current directory''. The final part names the script to execute.
The model uses Julia standard libraries; it does not require installing a
separate plotting package to generate its data.

The demonstration first prints the probabilities for the five-individual
example. They should be \(0.3\) for an increase, \(0.2\) for a decrease,
and \(0.5\) for no change. It then shows a short specified sequence of
draws, followed by a seeded repeated experiment.

The specified draws are a teaching example. They are not a random sample
used to estimate fixation probability. The repeated experiment that follows
uses the random-number generator.

\section{Read every part of the demonstration script}

The following blocks are inserted directly from \code{scripts/demo.jl}.
They are not separately maintained copies of the code. The numbers at the
left are the line numbers in the original file.

\subsection{Find the model file and load the libraries}

\lstinputlisting[firstline=1,lastline=5,firstnumber=1]{code/demo.jl}

Line 1 is a comment. The initial hash sign tells Julia that the rest of the
line is an explanation for people.

Line 2 loads the model file. The macro \code{@__DIR__} means the directory
containing this script. The function \code{joinpath} joins directory and
file names using the operating system's path convention. The component
\code{..} means ``go to the parent directory''. Starting from the
\code{scripts} folder, the constructed path goes up to the project folder,
then into \code{src}, then to \code{MoranModel.jl}. The function
\code{include} evaluates that file so its module becomes available.

Line 3 makes the exported names from the local \code{MoranModel} module
available in this script. The initial dot in \code{.MoranModel} identifies
a module already loaded in the current context. Line 4 loads the random-number
tools. Line 5 loads tools for formatted printing. The Julia documentation
describes these language and standard-library features~\cite{juliadocs}.

\subsection{Set the toy state and follow specified draws}

\lstinputlisting[firstline=7,lastline=20,firstnumber=7]{code/demo.jl}

Line 7 begins a function called \code{main}. This keeps the script's work
inside one function. Lines 8--10 assign the population size, current A count,
and selection advantage. Line 11 passes those three values to the model's
transition-probability function.

The returned value contains named fields. Thus \code{probabilities.up}
selects its increase probability, \code{probabilities.down} its decrease
probability, and \code{probabilities.stay} its no-change probability.
Lines 12--13 print those values for a reader to inspect.

The macro \code{@printf} prints values according to a format string.
In that string, \verb|%.3f| means a decimal number with three digits after
the decimal point. The character sequence \verb|\n| means a new line.
Formatting affects the displayed number, not the underlying calculation.

Line 15 creates an array: an ordered collection of six chosen draws. Square
brackets enclose its entries. Line 16 uses \code{enumerate} to obtain two
values for each entry: its position and the draw itself. The variable
\code{event} receives the position, and \code{u} receives the draw.

Line 17 recomputes the probabilities using the \emph{current} count.
Line 18 chooses the next count and assigns it back to \code{i}. This
assignment is what allows the count to evolve. Line 19 prints the event
number, draw, and new count. In its format, \verb|%d| means an integer,
while \verb|%.2f| means two digits after the decimal point. Line 20 ends
the loop.

\begin{workedexample}
At the first hand-trace event, the count is two and the increase threshold
is \(0.3\). The first specified draw is \(0.12\), which is smaller than
\(0.3\), so the count increases to three. Before the next draw, line 17
calculates new probabilities using three A individuals. The program does
not keep using the original thresholds from the two-individual state.
\end{workedexample}

\subsection{Repeat complete runs and summarise them}

\lstinputlisting[firstline=22,lastline=37,firstnumber=22]{code/demo.jl}

Line 23 creates a random-number generator with a recorded seed.
\code{MersenneTwister} is the generator's name; the integer in parentheses
initialises it. The variable \code{rng} stores the generator. Line 24
sets the number of independent simulation replicates to two thousand.
Line 25 starts the fixation counter at zero.

Line 26 repeats the following block once per replicate. Line 27 starts each
run at population size one hundred, one A individual, and advantage
\(s=0.01\). Each run continues until loss or fixation. The next replicate
starts again from one A individual. It does not inherit the preceding
replicate's final population.

The result has a Boolean field called \code{fixed}. A \textbf{Boolean}
is either \code{true} or \code{false}. In this addition, Julia treats
\code{true} as one and \code{false} as zero. Therefore
\code{fixations += outcome.fixed} adds one for a fixation and adds zero
for a loss. The operation \code{+=} means ``add to the existing value and
store the result back''.

After the loop ends, line 30 calculates the Wilson interval. Lines 31--34
print the fixation count, its estimated probability, its interval, and the
exact mathematical probability. The format \verb|%.6f| uses six digits
after the decimal point. The sequence \verb|%%| prints a literal percent
sign within an \code{@printf} format string.

Line 35 ends the definition of \code{main}. Line 37 calls it. Defining a
function describes its instructions; calling the function carries them out.

\section{Run the automated checks}

After the demonstration, run:
\begin{lstlisting}[language={},numbers=none]
julia --project=. test/runtests.jl
\end{lstlisting}
The checks ask whether the code obeys known properties of the model. They
cover valid inputs, boundary states, transition probabilities, exact fixation
probabilities, the independent backward-equation calculation, and other
specified behaviour. They do not prove that every real population follows
the Moran model. They assess whether this implementation does the things
we say it does.

Read a failure message before changing anything. A failing expected-value
check and a command that cannot locate a file are different problems. The
former may indicate incorrect code or an incorrect test. The latter means
the requested file was not found.

\section{Choose a quick study or the full study}

For a faster practice run, use:
\begin{lstlisting}[language={},numbers=none]
julia --project=. scripts/run_study.jl --quick
\end{lstlisting}
This runs one thousand replicates per parameter combination and saves data
under \code{results/demo}. It exercises the complete data-generation
procedure with a smaller Monte Carlo sample.

For the study used in this guide, omit the flag:
\begin{lstlisting}[language={},numbers=none]
julia --project=. scripts/run_study.jl
\end{lstlisting}
The default is twenty thousand replicates per combination. The full data
go under \code{results/full}. There is no \code{--full} flag; full mode
is simply the default. The script recognises \code{--quick} and rejects
unrecognised arguments rather than silently guessing what they mean.

The grid contains four population sizes and seven selection advantages,
giving \(4\times7=28\) combinations. We call one combination a
\textbf{cell}. The full run therefore contains \(28\times20000=560000\)
complete fixation experiments. The quick run contains
\(28\times1000=28000\). The deterministic calculations and small set of
display trajectories are generated in both modes.

Running a mode again replaces that mode's generated data files. If you want
to retain an experimental variant for comparison, first give the copied
project or its copied result folder a descriptive name.

\section{Read the study script: loading and writing a CSV row}

\lstinputlisting[firstline=1,lastline=12,firstnumber=1]{code/run_study.jl}

The loading instructions match the demonstration. This script also loads
\code{Dates}, which provides the finishing timestamp.

The helper function on lines 10--12 writes one row of values. Its argument
\code{file} identifies the open output file. Its argument \code{values}
contains the entries to write. The function \code{join(values, ",")}
turns those entries into text separated by commas. The function
\code{println(file, ...)} writes that text to the file and ends the line.

The resulting format is \textbf{CSV}, meaning comma-separated values.
Each line is a row; each comma separates columns. A spreadsheet program
can read these files, and they can also be opened as plain text. Our rows
contain numbers or simple labels that do not themselves contain commas.
The helper is sufficient for these controlled outputs; it is not a general
CSV parser for arbitrary quoted text.

\section{Read the study script: the fixation grid}

\lstinputlisting[firstline=15,lastline=44,firstnumber=15]{code/run_study.jl}

\subsection{Open the file and label its columns}

Line 15 defines a function taking three inputs: where to save results,
how many replicates to run, and a base seed. Line 16 constructs the name
\code{fixation.csv}. Line 17 opens that file for writing. The string
\code{"w"} specifies write mode.

The syntax \code{do file} starts a \textbf{do block}. It gives the opened
file the local name \code{file} within the block. Julia closes the file
when the block finishes, including when an error interrupts the block.
Line 18 writes the header: the names of the data columns.

\subsection{Visit all parameter combinations}

Line 19 starts the cell counter at zero. Line 20 selects a population size
from an array. Inside that loop, line 21 selects each advantage from a second
array. These are \textbf{nested loops}: all seven advantages are used for
one population size before moving to the next size.

Line 22 increments the cell number. Line 23 makes its seed by adding the
cell number to the base seed. Line 24 constructs the generator for this
cell. This gives every cell a recorded seed without resetting the generator
at the start of every replicate. The generator continues its sequence as
replicates run.

\subsection{Run each replicate to an absorbing state}

Lines 25--26 initialise two counters. One counts fixations; the other sums
the number of event attempts used by the complete runs. Lines 27--31 repeat
the experiment. The initial A count in line 28 is always one.
Line 29 adds the Boolean fixation result. Line 30 adds that run's event
count, including unchanged attempts.

There is no display horizon in this loop. It continues each replicate until
one of the two absorbing boundaries is reached. Thus the fixation estimate
is not contaminated by treating an unfinished display trajectory as a loss.

\subsection{Summarise and save the cell}

Line 32 divides fixations by replicates. Line 33 calculates the interval.
Line 34 evaluates the exact formula, independently of the simulation count.
Line 35 computes the usual plug-in standard-error estimate for a proportion:
the square root of the estimated probability times one minus that estimate,
divided by the number of replicates. This standard error describes the
approximate Monte Carlo sampling scale. It is not another biological parameter.
The plotted confidence interval uses the Wilson calculation instead.

Lines 36--38 pass one \textbf{tuple} of results to the row writer. A tuple
is an ordered group of values, written here in parentheses. The order must
match the header. In order, the columns are:
\begin{itemize}
 \item \code{N}, \code{s}, and \code{i0}: population size, advantage,
       and initial A count;
 \item \code{repetitions} and \code{fixations}: replicate count and
       observed fixation count;
 \item \code{estimate}, \code{lower}, and \code{upper}: estimated
       probability and Wilson interval endpoints;
 \item \code{exact} and \code{standard_error}: the mathematical target
       and estimated Monte Carlo standard error;
 \item \code{events_mean} and \code{seed}: average number of event
       attempts to either absorbing outcome, and the cell's seed.
\end{itemize}
The average event count includes both fixation and loss runs. It is not
the average time conditional on successful fixation.

Line 39 uses \code{flush} to send buffered output to the file during the
study. Line 40 prints a progress update. In its format, \verb|%2d| reserves
at least two character positions for an integer and \verb|%3d| reserves
at least three. The widths improve alignment; they do not limit the numbers.
The remaining \code{end} statements finish the inner loop, outer loop,
file block, and function.

\section{Read the study script: compare initial conditions}

\lstinputlisting[firstline=47,lastline=59,firstnumber=47]{code/run_study.jl}

This block opens \code{initial_frequency.csv} and loops over the same
population sizes and advantages. It calculates probabilities with the exact
formula, so there is no random sampling in this block.

Line 52 uses Julia's \code{div(N, 10)} function for integer division,
to calculate \code{N} divided by ten as an integer count. All four supplied
population sizes are divisible by ten. Thus the result is exactly ten percent
of the population: two, five, ten, or twenty A individuals.

Line 53 calculates fixation probability from one copy. Line 54 calculates
it from the ten-percent count. The row in line 55 records both answers,
the ten-percent initial count and frequency, and the single-copy fixation
probability relative to its neutral baseline.

For the last comparison, the neutral single-copy probability is \(1/N\).
Dividing the selected probability by \(1/N\) is the same as multiplying
it by \(N\), which explains \code{N * single}. A value of two in that
column means twice the neutral odds. It does not mean a probability of two.

\section{Read the study script: check the approximation}

\lstinputlisting[firstline=62,lastline=75,firstnumber=62]{code/run_study.jl}

This block writes \code{approximation.csv}. Line 65 creates 121 positive
selection values from \(10^{-5}\), which is \(0.00001\), up to \(0.5\).
They are evenly spaced on a logarithmic scale, not evenly spaced by ordinary
addition. That gives detail near very small advantages as well as larger ones.

The function \code{log10} is the logarithm to base ten: it tells us the
power to which ten must be raised. The function \code{range} produces
evenly spaced exponents between the two endpoint logarithms. Its
\code{length=121} argument specifies the number of values. The operator
\code{.^} raises ten to \emph{each} exponent in that range. The dot
means to apply the operation element by element; this is called
\textbf{broadcasting}.

Lines 66--67 loop over sizes and these finer selection values. Line 68
calculates the exact single-copy answer. Line 69 calculates the
weak-selection approximation explained earlier. The use of
\code{-expm1(-N * s)} calculates \(1-\exp(-Ns)\) accurately near zero.
Line 70 subtracts the exact answer from the approximation and divides by
the exact answer. This is a \emph{signed} relative error: positive means
the approximation is above the exact answer, and negative means it is below.

The values in this block are mathematical comparisons, not estimated
simulation frequencies. Increasing the number of replicates cannot repair
an inaccurate approximation. The loop starts above zero; the neutral limit
is explained separately instead of evaluating a zero-over-zero expression.

\section{Read the study script: save example trajectories}

\lstinputlisting[firstline=78,lastline=119,firstnumber=78]{code/run_study.jl}

\subsection{Set the display experiment}

Lines 79--82 choose a population of one hundred, a display horizon of thirty
thousand event attempts, output every one hundred attempts, and eight paths
per condition. A \textbf{path}, or trajectory, is the sequence of counts in
one run. Line 83 starts a condition counter. Lines 84--86 open the trajectory
summary file and label its columns.

Lines 87--88 combine initial counts of one and ten with three advantages:
zero, \(0.01\), and \(0.1\). There are six conditions in total. Line 89
updates the condition number.

\subsection{Store eight paths without selecting their outcomes}

Line 90 creates an initially empty vector whose entries will themselves be
vectors of integer counts. The notation \code{Vector{Vector{Int}}()} means
``a vector of integer vectors''. Line 91 repeats the path-generation block
eight times.

Line 92 assigns a seed using the base seed, the condition number, and the
path number. The expression \code{100condition} means
\code{100 * condition}; Julia allows a numerical coefficient directly
before a variable name. The large fixed offset keeps these path seeds
separate from the fixation-grid seeds. Line 93 creates a generator with that
seed and simulates the path. Line 94 appends its count vector to the
collection. The exclamation mark in \code{push!} indicates that it changes
the collection it receives.

These are predetermined seeds and path counts. The script does not search
for an appealing successful trajectory and discard the others.

Lines 95--98 inspect when the trajectory first reached a boundary.
The value \code{nothing} means no absorption event was recorded before
the display horizon. The comparison \code{=== nothing} tests for that
particular missing-value marker. For the CSV file, the code replaces it
with \(-1\). This negative number is a label for ``not yet absorbed within
the display'', not a negative event time and not a loss.

Lines 99--100 write the condition, path number, seed, absorption indicator,
absorption time or marker, and final count. The expression
\code{trajectory.counts[end]} retrieves the last entry of the counts vector.
The \code{end} inside square brackets means ``the final index'' here; it
is not the instruction that closes a loop.

\subsection{Save frequencies on a readable time grid}

Line 102 calculates the deterministic comparison for the same starting
state and horizon. Line 103 turns the selection value into text and replaces
its decimal point with the letter \code{p}. In \code{"." => "p"}, the
arrow specifies a replacement rule. Thus \code{0.01} becomes the filename
label \code{0p01}.

Line 104 inserts the parameter values into a filename. Inside a Julia
string, \code{$(N)} means ``evaluate this expression and insert its text
here''. This is \textbf{string interpolation}. For example, one generated
name is \code{trajectory_N100_i1_s0p01.csv}.

Lines 105--106 open the path file and write its header. Line 107 uses this range:
\begin{lstlisting}[numbers=none]
0:sample_every:max_events
\end{lstlisting}
Its three components mean starting
value, step size, and final allowed value. Here it visits events zero, one
hundred, two hundred, and so on, through thirty thousand.

Line 108 creates the beginning of a numerical row. The type
\code{Float64} requests ordinary double-precision numerical entries. The
first entry is the event number; the second is the event number divided by
the population size. That second number is our generation-like time unit.
At one hundred events and \(N=100\), it equals one. No synchronous
replacement of all individuals is implied.

Lines 109--111 append each path's frequency at that event. The array index
is \code{event + 1} because array position one contains the initial state
at event zero. Dividing the stored count by \code{N} converts it into a
frequency. Line 112 appends the deterministic frequency, which is already
a frequency and must not be divided by \code{N} again. Line 113 writes
the row. The remaining \code{end} statements close the loops, file blocks,
and function.

The stored histories extend to the display horizon. After absorption, a
history repeats its absorbing count. That is why a path can stay exactly at
zero or one for the rest of the displayed time.

\section{Read the study script: record the run information}

\lstinputlisting[firstline=122,lastline=142,firstnumber=122]{code/run_study.jl}

This function writes \code{metadata.txt}, a plain-text description of the
computation. The file records the Julia version, random-number generator,
base seed, seed conventions, parameter grid, replicate count, and mode.
The constant \code{VERSION} contains the running Julia version.

The file also records the stopping rules. This matters because the fixation
experiment runs until absorption, whereas displayed trajectories have a
finite horizon. A future reader should not have to guess which rule was used.

Line 139 writes the elapsed time supplied by the main function. Line 140
writes the completion timestamp in UTC using \code{Dates.now(Dates.UTC)}.
A timestamp identifies when the computation finished; elapsed seconds measure
how long it took. They serve different purposes. Runtime depends on the
computer and Julia compilation state, so it is not a biological result.

\section{Read the study script: the main function}

\lstinputlisting[firstline=145,lastline=168,firstnumber=145]{code/run_study.jl}

Line 146 checks whether the string \code{"--quick"} is present in
\code{ARGS}, Julia's list of command-line arguments supplied after the
script name. Line 147 uses \code{setdiff} to identify supplied arguments
outside the permitted list. Line 148 checks whether that unrecognised list
is empty.

The operator \code{||} means short-circuit ``or''. If the left expression
is true, Julia does not evaluate the right expression. Thus an empty
unrecognised-argument list allows execution to continue; otherwise
\code{error} reports which arguments were not understood. This is an input
check, not a statistical test.

Lines 149--150 set the default replicate count and output mode. Lines
151--154 replace them with one thousand and \code{"demo"} when quick
mode was requested. Line 155 sets the base seed. Line 156 constructs the
output directory relative to the script's location; \code{normpath}
simplifies path components such as \code{..}. Line 157 uses \code{mkpath}
to create the directory if needed.

Line 158 records the clock time before the work begins. Lines 159--162 call
the four data-generation functions in order. Line 163 subtracts the starting
clock time from the current clock time, producing elapsed seconds. Line 164
writes the metadata. Line 165 prints a completion message; its \verb|%s|
placeholder prints text. Line 166 ends the function definition, and line
168 calls the function to run the script.

\section{Turn full results into the document inputs}

After generating the full results, run:
\begin{lstlisting}[language={},numbers=none]
julia --project=. scripts/prepare_guide.jl
\end{lstlisting}
This preparation step reads the full-study CSV files and writes the numerical
and figure inputs used by the LaTeX document. It does not run a new set of
fixation simulations. This separation ensures that the written results and
figures refer to the saved study data.

The guide uses the full results, not the smaller quick demonstration. If
you rerun only quick mode, you have practised the pipeline but have not
replaced the full tables used in the guide. To update the guide with a
changed full experiment, rerun full mode and this preparation step.

\section{Compile or open the LaTeX guide}

If Tectonic is available, run these commands from the project folder:
\begin{lstlisting}[language={},numbers=none]
cd guide
tectonic --keep-logs main.tex
\end{lstlisting}
The first command moves into the guide folder. The argument \code{main.tex}
identifies the main LaTeX source file. Tectonic writes \code{main.pdf} in that
folder. The option keeps a compilation log that can help diagnose errors.

You can also open the \code{guide} project folder in the Codex LaTeX
editor and use its PDF preview. The LaTeX compilation step uses the saved
source and prepared inputs. It does not execute Julia automatically.

An online editor such as Overleaf can be used for editing LaTeX too. The
preparation helper makes the guide folder self-contained: it copies the actual
runnable source into \code{guide/code} and the plotting data into
\code{guide/generated}. Upload the contents of the whole guide folder, choose
\code{main.tex} as the main document, and use pdfLaTeX. Uploading only the main
file would omit the chapters and supporting inputs. Edit the original Julia
files and rerun the helper to refresh their copies before recompiling.

\section{Make a small change and understand what it changes}

Keep an unchanged copy of the supplied project as a reference. In a working
copy, change one setting at a time. First make a prediction in words. Then
run the demonstration or quick study. Finally compare the observation with
your prediction and the exact result.

For example, change only the demonstration's selection advantage from
\(0.5\) to zero in the toy-state block. Keep \(N=5\) and \(i=2\).
Predict that the increase and decrease probabilities will become equal.
Their new values should be \(0.24\), while the no-change probability
should be \(0.52\). This is a stronger learning exercise than changing
several parameters and guessing which one caused the difference.

Changing a printed label does not change a function's input. Check the
actual arguments passed to the function. In particular, the seeded experiment
in the demonstration passes its own numbers explicitly, so changing the
earlier toy variables does not change that later experiment.

If you add a new population size to the study, remember that the exact
ten-percent comparison currently uses integer division by ten. Choose a
size divisible by ten if you want to preserve that initial frequency exactly.
Changes to the scientific question may require changes to more than one
output label or loop; work through them deliberately.

\section{Diagnose a message before drawing a conclusion}

\begin{itemize}
 \item \textbf{The terminal says Julia was not found.} The operating system
       could not locate the Julia executable. Check the Julia installation
       and terminal path. No simulation result has been produced.
 \item \textbf{A script or source file cannot be opened.} Check the current
       folder and the extracted directory layout. Keep the project folders
       together.
 \item \textbf{An \code{ArgumentError} describes an input.} Read that
       description. The model requires an integer population size of at
       least two, an integer count between zero and the size, and a finite
       nonnegative advantage. A rejected input is not evidence about allele
       behaviour.
 \item \textbf{An unknown command-line argument is reported.} Use no flag
       for full mode or \code{--quick} for the demonstration mode.
 \item \textbf{A simulation estimate differs from the exact probability.}
       Compare the difference with its sampling uncertainty. A small
       discrepancy is expected. Repeated systematic discrepancies much
       larger than that scale deserve investigation.
 \item \textbf{A display path has absorption time \(-1\).} It had not
       reached a boundary by the displayed horizon. It was not automatically
       counted as a loss in the separate fixation experiment.
\end{itemize}

\begin{checkpoint}
\textbf{Question 1.} Which command produces the twenty-thousand-replicate
study?

\textbf{Answer.} Run \code{julia --project=. scripts/run_study.jl} from
the project folder, without \code{--quick}.

\textbf{Question 2.} You edited the full-study settings, ran only quick
mode, and recompiled the guide. Should its full-study tables have changed?

\textbf{Answer.} No. Quick mode writes the demonstration results. Rerun
the full study, run \code{scripts/prepare_guide.jl}, and then compile the
document to update the full-study presentation.

\textbf{Question 3.} A trajectory file reports event \(500\) at generation
\(5\) for \(N=100\). What calculation produced that time?

\textbf{Answer.} The script divided the event count by the population size:
\(500/100=5\). It is a model time unit, not five synchronised replacements
of every individual.
\end{checkpoint}
